\chapter{Un mini terminal, multiplateforme}

Un terminal semble simple : on tape une commande, on lit la réponse. C'est
trompeur. Un vrai terminal doit lancer un programme, lui parler pendant qu'il
tourne, comprendre un langage de contrôle vieux de cinquante ans, et ne jamais
figer l'interface --- sur trois systèmes qui ne s'y prennent pas du tout de la
même façon.

NKCode en contient un, complet. Il tient en trois fichiers, et ce découpage est
tout le chapitre.

\begin{center}
\begin{tabular}{@{}llr@{}}
\toprule
\textbf{Fichier} & \textbf{Rôle} & \textbf{Lignes} \\
\midrule
\texttt{NkProcess.h} & lancer une commande, lire sa sortie & 171 \\
\texttt{NkPty.h/.cpp} & tenir un shell \emph{interactif} vivant & 496 \\
\texttt{NkTerm.h} & interpréter le flux, en faire une grille & 631 \\
\bottomrule
\end{tabular}
\end{center}

\section{Niveau 1 : lancer une commande et ne pas geler}

\begin{nkcode}{Applications/NKCode/src/NKCode/Project/NkProcess.h}
//   Start(cmd) lance la commande sur un THREAD d'arriere-plan (stdout+stderr
//   fusionnes) ; le thread pousse chaque ligne dans une file THREAD-SAFE ;
//   l'UI appelle Drain() chaque frame pour recuperer les nouvelles lignes sans
//   geler.
\end{nkcode}

\begin{nkcode}{Applications/NKCode/src/NKCode/Project/NkProcess.h}
void Drain(NkVector<NkString> &out) {
    threading::NkScopedLock<NkMutex> lk(mMutex);
    for (/* ... */)
        out.PushBack(mLines[i]);
    // ...
}
\end{nkcode}

\begin{nkretenir}
\textbf{C'est le patron « producteur lent, interface fluide », et il vaut bien
au-delà des terminaux.}

Un fil d'arrière-plan lit le programme et pousse ses lignes dans une file
protégée par un mutex. L'interface, à chaque image, appelle \texttt{Drain} : elle
prend ce qui est arrivé, ou rien, et rend la main immédiatement.

Elle n'\emph{attend} jamais. Une compilation de trois minutes défile ligne à ligne
pendant que la fenêtre reste vivante.
\end{nkretenir}

\begin{nkpiege}
\textbf{Le verrou est pris pendant la copie, pas pendant la lecture du
programme.} C'est la règle du chapitre NKThreading appliquée : \emph{verrouiller
le temps de toucher la donnée, pas le temps de travailler}.

Si le fil gardait le verrou pendant qu'il attend la ligne suivante --- ce qui peut
durer des secondes --- l'interface bloquerait à chaque image en essayant de le
prendre. Vous auriez payé un fil pour obtenir un programme séquentiel qui
saccade.
\end{nkpiege}

\begin{nkretenir}[Une exception à la règle, et elle est écrite]
\texttt{NkProcess.h} inclut \texttt{<cstdio>}, alors que le dépôt interdit la
bibliothèque standard. Son en-tête le dit et le justifie :

\begin{quote}
« WRAPPER MAISON DÉSIGNÉ pour les PIPES process : les appels
\texttt{\_popen}/\texttt{fgets}/\texttt{fgetc} de bas niveau sont volontairement
conservés dans cette famille de wrappers (pas d'équivalent \texttt{NkFile} pour un
flux de processus). »
\end{quote}

Comparez au chapitre précédent, où \texttt{NkVideoPlayer} incluait les mêmes
en-têtes \emph{sans rien dire}. Le code est identique ; ce qui diffère est
décisif.

\textbf{Une exception écrite est une décision --- on peut la discuter, la lever,
la dater. Une exception muette est une dette que le prochain lecteur prendra pour
un exemple à suivre.}
\end{nkretenir}

\section{Niveau 2 : un shell qui reste vivant}

\texttt{NkProcess} lance une commande, la lit, et meurt. Ça ne suffit pas : un
terminal doit garder \emph{le même} shell entre deux commandes --- sinon
\texttt{cd} ne change rien, et les variables s'oublient.

\begin{nkcode}{Applications/NKCode/src/NKCode/Project/NkPty.h}
// Contrairement a NkProcess (un _popen jetable par commande), NkPty lance UN
// shell (powershell / wsl -d <distro> / bash / cmd) attache a un pseudo-console
// Windows (ConPTY). On ECRIT les frappes clavier dans son entree et on LIT en
// continu sa sortie (avec sequences VT : couleurs, deplacement curseur, clear).
// => le shell affiche lui-meme son invite ; `clear` efface ; pas de boite de
//    saisie separee.
\end{nkcode}

\begin{nkretenir}
\textbf{« Le shell affiche lui-même son invite » est le renversement complet.}

Le débutant écrit une zone de saisie, un bouton \emph{Exécuter}, et une zone de
sortie. C'est une \emph{imitation} de terminal : l'invite est fausse, les couleurs
sont perdues, \texttt{vim} ne marche pas, \texttt{clear} n'efface rien.

Un vrai terminal ne fabrique rien. Il transmet les frappes au shell et affiche ce
que le shell renvoie --- \emph{y compris son invite}. Il n'est pas l'auteur de ce
qu'on voit ; il en est le verre.
\end{nkretenir}

\section{Le pseudo-terminal : ce que c'est, et pourquoi il en faut un}

Un programme se comporte différemment selon qu'il écrit dans un fichier ou vers un
humain : vers un fichier, il retire les couleurs et l'interactivité. Un tube
ordinaire ressemble à un fichier --- d'où des sorties ternes et \texttt{vim} qui
refuse de démarrer.

Un \textbf{pseudo-terminal} est un tube qui se \emph{présente} comme un terminal :
il a une taille en lignes et colonnes, il accepte d'être redimensionné, et le
programme d'en face le croit.

\begin{nkcode}{Applications/NKCode/src/NKCode/Project/NkPty.h}
bool Start(const NkString &cmdline, int16 cols, int16 rows,
           const NkString &cwd = NkString());
void Write(const char *data, usize len);   // les frappes clavier
void Resize(int16 cols, int16 rows);       // quand le panneau change de taille
\end{nkcode}

\begin{nkpiege}
\textbf{\texttt{Resize} n'est pas cosmétique.} Le programme d'en face décide sa
mise en page d'après la taille qu'on lui a déclarée. Oubliez de l'informer, et
\texttt{htop} dessinera pour 80 colonnes dans une fenêtre qui en fait 200 --- ou
débordera dans l'autre sens.

Un terminal qui ne réagit pas au redimensionnement n'est pas « moins abouti » : il
est faux dès la première fenêtre de taille inhabituelle.
\end{nkpiege}

\section{Le multiplateforme, et où il est isolé}

\begin{nkcode}{Applications/NKCode/src/NKCode/Project/NkPty.cpp (structure)}
#if defined(_WIN32)
    // ConPTY : CreatePseudoConsole, PROC_THREAD_ATTRIBUTE_PSEUDOCONSOLE
#else
  #if defined(__APPLE__)
    // <util.h>
  #else
    // <pty.h>   -- forkpty
  #endif
#endif
\end{nkcode}

\begin{nkretenir}
\textbf{Trois systèmes, trois mécanismes, un seul en-tête.} \texttt{NkPty.h}
n'expose que \texttt{Start}, \texttt{Write}, \texttt{Read}, \texttt{Resize} : rien
de ce qui suit ne connaît ConPTY ni \texttt{forkpty}.

L'en-tête dit d'ailleurs \emph{pourquoi} le \texttt{.cpp} existe : « isolée dans
\texttt{NkPty.cpp} pour ne pas polluer les en-têtes GUI avec les macros de
\texttt{windows.h} ». Ce n'est pas une préférence de style --- \texttt{windows.h}
définit des macros comme \texttt{min}, \texttt{max} et \texttt{None} qui cassent
le code qui les inclut par accident.

\textbf{Là où le portage est le plus sale, l'interface doit être la plus
étroite.}
\end{nkretenir}

\section{Niveau 3 : comprendre ce que le shell renvoie}

Le shell ne renvoie pas du texte : il renvoie du texte \emph{mêlé d'ordres}.
\texttt{ESC[31m} veut dire « à partir d'ici, écris en rouge », \texttt{ESC[2J}
« efface l'écran », \texttt{ESC[10;5H} « place le curseur ligne 10, colonne 5 ».

\begin{nkcode}{Applications/NKCode/src/NKCode/Project/NkTerm.h}
// NkTerm.h - Emulateur de terminal VT : transforme le flux d'octets d'un shell
//   en une GRILLE de cellules (caractere + couleur avant/arriere) avec curseur,
//   defilement et scrollback.
//   Rendu : l'UI lit TotalLines()/LineAt() et dessine chaque cellule (monospace).

struct NkTermCell {
        uint32 cp = 0x20;                  // codepoint Unicode
        NkColor fg = {204, 204, 204, 255};
        NkColor bg = {0, 0, 0, 0};         // a==0 => fond par defaut
};
\end{nkcode}

\begin{nkretenir}
\textbf{L'émulateur transforme un flux en état.} Il consomme des octets et
maintient une grille : à chaque case, un caractère et deux couleurs. Le dessin,
ensuite, est trivial --- une boucle sur des cellules, en police à chasse fixe.

\textbf{C'est la séparation qui rend l'affaire faisable} : d'un côté un automate
qui ne dessine rien et se teste sans fenêtre ; de l'autre un dessin qui ne
comprend rien et ne peut donc pas se tromper. Exactement le découpage
règles/écran des jeux, appliqué à un terminal.
\end{nkretenir}

\begin{nkpiege}
\textbf{Une séquence d'échappement peut arriver coupée en deux.} La lecture rend
ce qui est disponible : \texttt{ESC[3} dans un paquet, \texttt{1m} dans le
suivant.

Un analyseur qui traite chaque paquet indépendamment affichera \texttt{[3} puis
\texttt{1m} en clair, aléatoirement, une fois de temps en temps --- le défaut qui
ne se reproduit jamais chez le développeur et toujours chez l'utilisateur.

\textbf{L'automate doit garder son état d'un paquet à l'autre.} C'est le même
problème que le commentaire de bloc du chapitre précédent : \emph{un analyseur
travaille sur un flux, pas sur des morceaux}.
\end{nkpiege}

\begin{nkretenir}
Notez \texttt{bg = \{0,0,0,0\}} avec le commentaire « \texttt{a==0} $\Rightarrow$
fond par défaut ». Une cellule ne dit pas « fond noir » : elle dit « je n'ai pas
d'avis ».

C'est ce qui permet au terminal de suivre le thème de l'éditeur. Une valeur
\emph{absente} et une valeur \emph{noire} sont deux choses différentes --- et les
confondre est l'erreur qui rend une interface non thémable.
\end{nkretenir}

\begin{nkbilan}
Un terminal se construit en trois étages qui se justifient l'un l'autre :
\emph{lancer} sans geler (fil d'arrière-plan, file protégée, \texttt{Drain} par
image), \emph{tenir un shell vivant} via un pseudo-terminal --- c'est le shell qui
affiche son invite, pas nous ---, et \emph{interpréter} le flux VT en une grille
de cellules. Le code spécifique à chaque système est isolé dans un seul
\texttt{.cpp}, derrière quatre fonctions. Et l'automate garde son état entre les
paquets, sinon les séquences coupées s'affichent en clair, rarement.
\end{nkbilan}

\begin{nkquestions}
\begin{enumerate}
\item Pourquoi \texttt{Drain} plutôt qu'attendre la fin de la commande ?
\item Pourquoi un tube ordinaire ne suffit-il pas pour un shell interactif ?
\item Que se passe-t-il si l'on n'appelle jamais \texttt{Resize} ?
\item Pourquoi le code Win32 est-il dans un \texttt{.cpp} et non dans l'en-tête ?
\item Pourquoi une cellule a-t-elle un fond d'alpha nul plutôt qu'un fond noir ?
\end{enumerate}
\end{nkquestions}

\begin{nkpratique}
Construisez votre mini terminal, en trois étapes qui correspondent aux trois
niveaux du chapitre.

\textbf{A.} Une zone de saisie, un bouton, une zone de sortie. Lancez la commande
sur un fil, videz la file à chaque image. \emph{Vérifiez que l'interface reste
fluide} pendant une commande de trente secondes.

\textbf{B.} Remplacez par un shell persistant. \texttt{cd} doit désormais avoir un
effet sur la commande suivante --- c'est votre critère de réussite, et il est
binaire.

\textbf{C.} Interprétez au minimum les couleurs (\texttt{SGR}) et l'effacement
d'écran (\texttt{ED}). Le reste des séquences peut être \emph{consommé et ignoré}
--- mais consommé, jamais affiché.
\end{nkpratique}

\begin{nkdemo}
Prouvez que votre analyseur survit à une séquence coupée.

Écrivez un banc \emph{sans fenêtre} qui alimente votre émulateur avec
\texttt{"ESC[31mROUGE"} en un seul morceau, puis avec le même texte découpé
\textbf{à chaque octet} --- un appel par caractère. Comparez la grille obtenue,
cellule par cellule.

Les deux doivent être identiques. Faites ensuite varier le point de coupure sur
toutes les positions possibles et vérifiez les $n$ résultats.

\textbf{Pourquoi la variation exhaustive} : une coupure à un seul endroit peut
tomber juste par chance. C'est la coupure \emph{au milieu du nombre} ---
\texttt{ESC[3} puis \texttt{1m} --- qui casse les analyseurs naïfs, et rien ne
garantit que votre essai unique tombe dessus.
\end{nkdemo}

\begin{nklogique}
Votre terminal affiche parfois \texttt{[0m} en clair au milieu de la sortie ---
environ une fois sur cent lancements, jamais au même endroit.

\begin{enumerate}
\item Énumérez au moins trois causes possibles.
\item Pour chacune, dites quelle mesure la distinguerait des autres --- une mesure
      qui donne un résultat \emph{différent} selon la cause, sinon elle ne tranche
      rien.
\item Un collègue propose de filtrer \texttt{[0m} avant l'affichage. Expliquez
      pourquoi ce correctif est plus dangereux que le défaut, en une phrase.
\end{enumerate}
\end{nklogique}
